\section{Билет 41. Степени свободы движения частицы вещества. Закон равнораспределения энергии по степеням свободы. Средняя кинетическая энергия поступательного движения одноатомной молекулы и ее связь с температурой. Число степеней свободы и средняя энергия многоатомной молекулы.}

\begin{definition}
\textbf{Числом степеней свободы} $i$ механической системы называется минимальное число независимых переменных (координат), полностью определяющих положение этой системы в пространстве.
\begin{itemize}
    \item Для одноатомной молекулы (материальная точка) $i = 3$ (поступательные координаты $x, y, z$).
    \item Для двухатомной молекулы (жёсткая гантель) добавляются 2 вращательные степени свободы, итого $i = 5$ (3 поступательные + 2 вращательные).
    \item Для многоатомной молекулы ($N \ge 3$ атомов) добавляются 3 вращательные степени свободы, итого $i = 6$ (при отсутствии колебаний).
\end{itemize}
\end{definition}

\begin{theorem}[Закон равнораспределения энергии по степеням свободы]
При термодинамическом равновесии на каждую поступательную и вращательную степень свободы молекулы в среднем приходится одинаковая кинетическая энергия:
\begin{equation}
    \langle \epsilon_{1} \rangle = \frac{1}{2} k T, \label{eq:equipart_41}
\end{equation}
где $k$ --- постоянная Больцмана, $T$ --- абсолютная температура.
Если молекула обладает колебательными степенями свободы, то на каждую из них приходится энергия $kT$, так как колебательная степень свободы обладает как кинетической, так и потенциальной энергией (каждая в среднем равна $kT/2$).
\end{theorem}

\begin{property}[Средняя энергия одноатомной молекулы]
Одноатомная молекула имеет только три поступательные степени свободы ($i=3$). Согласно закону равнораспределения, её средняя кинетическая энергия поступательного движения равна:
\begin{equation}
    \langle E_{\text{пост}} \rangle = 3 \cdot \frac{1}{2} k T = \frac{3}{2} k T. \label{eq:mono_energy_41}
\end{equation}
Эта формула устанавливает прямой статистический смысл абсолютной температуры: температура является мерой средней кинетической энергии хаотического поступательного движения молекул идеального газа.
\end{property}

\begin{definition}[Многоатомные молекулы]
Для молекулы, состоящей из $N$ атомов, общее число степеней свободы складывается из поступательных, вращательных и колебательных:
\begin{equation}
    i = i_{\text{пост}} + i_{\text{вр}} + 2i_{\text{кол}}.
\end{equation}
Число колебательных степеней свободы определяется как $i_{\text{кол}} = 3N - i_{\text{пост}} - i_{\text{вр}}$.
Средняя полная энергия (кинетическая + потенциальная) такой молекулы при тепловом равновесии равна:
\begin{equation}
    \langle \epsilon \rangle = \frac{i}{2} k T. \label{eq:poly_energy_41}
\end{equation}
При обычных температурах колебательные степени свободы «заморожены» (квантовый эффект), поэтому для двухатомных газов обычно берут $i=5$, а для многоатомных (нелинейных) $i=6$.
\end{definition}
